\subsection*{4. Regular polyhedra and finite fields}
The fourth section starts from the regular maps whose automorphisms act
transitively on frames, translates their oriented version into Galois
coverings by the dictionary of \S3, and includes infinite maps and
tessellations (pp.~18--19). It then follows the geometric realization of
vertices, edges, and faces into affine and projective spaces over arbitrary
base rings. The move to projective space on p.~20 supplies the author's
duality formalism and accommodates centers at infinity. ``Polyèdres
épinglés'', the higher cartographic group, fundamental reflections, and
parameters expressed through doubled cosines lead on p.~21 to the
universal polynomial base ring; specializations over finite fields then
lead back to the arithmetic curves of \S3 (pp.~22--23). A mathematical
slip occurs across pp.~21--22, where the group generated by the fundamental
reflections is identified with the full automorphism group, and the
icosahedron's group is then called simple of order 60. On the usual
centered icosahedron with vertices
$\{(0,\pm1,\pm\varphi),(\pm1,\pm\varphi,0),(\pm\varphi,0,\pm1)\}$,
$\varphi=(1+\sqrt5)/2$, the map $j:v\mapsto-v$ preserves vertices,
distances, and hence all incidences. It is a central involution in its
orthogonal symmetry group. The cyclic coordinate permutation
$(x,y,z)\mapsto(y,z,x)$ also preserves that vertex set and has order three,
so $\langle j\rangle$ is a nontrivial proper central subgroup of the full
group. The determinant morphism to $\{1,-1\}$ selects the
orientation-preserving subgroup as its kernel. The earlier projective
context leaves this calculation intact: on the affine chart of projective
three-space the same map is represented by
$\operatorname{diag}(-1,-1,-1,1)$ and sends $[v:1]$ to $[-v:1]$.
The readers preserve the scanned statement. On pp.~23--24 the
``paradigme combinatoire'' names a further direction suggested by
specialization in exceptional characteristics; the discussion also names
the 27 lines on a cubic surface and the root system $E_7$, which is
retained. Grothendieck records the limits of the work he has actually
pursued, introduces the circumscribed hyperquadric, and on p.~25 proposes
a modern rewriting of Klein's book as a task for a young researcher.
Coxeter's books and Klein's book are the works explicitly invoked here;
the geometric-combinatorial viewpoint continues into the foundational
questions of \S5.
(\href{https://grothendieck.umontpellier.fr/web/viewer.html?file=../119.pdf\#page=21}{Copy 1, pp.~18--25, beginning at site 20};
\href{https://grothendieck.umontpellier.fr/web/viewer.html?file=../119.pdf\#page=24}{reflection-group wording on pp.~21--22, site 23};
\href{https://grothendieck.umontpellier.fr/web/viewer.html?file=../119.pdf\#page=25}{site 24}.)
% Complete the reference to note (5) when its scan check is finished.

